% =====================================================================
% AcolheMente Escolar PB — Documento Acadêmico Final
% Classe: abnTeX2 | Estilo: abntex2-alf | Compilação: xelatex + bibtex
% =====================================================================
\documentclass[
    12pt,
    openright,
    oneside,
    a4paper,
    english,
    brazil
]{abntex2}

% --- Impede que "Sumário" apareça como item do próprio sumário ---
\renewcommand{\tocheadstart}{}
\renewcommand{\aftertoctitle}{}

% --- Pacotes obrigatórios ---
\usepackage[utf8]{inputenc}
\usepackage[T1]{fontenc}
\usepackage{indentfirst}
\usepackage[brazil]{babel}
\usepackage[alf]{abntex2cite}
\usepackage{float}
\usepackage{url}
\usepackage{booktabs}
\usepackage{graphicx}
\usepackage{xcolor}
\usepackage{listings}
\usepackage{longtable}
\usepackage{amsmath}
\usepackage{amssymb}
\usepackage[aboveskip=6pt,position=top]{caption}
\captionsetup[table]{position=top}
\captionsetup[figure]{position=bottom}

% --- Remoção de bordas vermelhas nos links ---
\hypersetup{
    colorlinks=false,
    pdfborder={0 0 0},
    hidelinks
}

% --- Configuração de listagens de código ---
\definecolor{codegreen}{rgb}{0,0.6,0}
\definecolor{codegray}{rgb}{0.5,0.5,0.5}
\definecolor{codepurple}{rgb}{0.58,0,0.82}
\definecolor{backcolour}{rgb}{0.95,0.95,0.92}
\lstset{
    backgroundcolor=\color{backcolour},
    commentstyle=\color{codegreen},
    keywordstyle=\color{magenta},
    numberstyle=\small\color{codegray},
    stringstyle=\color{codepurple},
    basicstyle=\ttfamily\small,
    breakatwhitespace=false,
    breaklines=true,
    captionpos=t,
    keepspaces=true,
    numbers=left,
    numbersep=5pt,
    showspaces=false,
    showstringspaces=false,
    showtabs=false,
    tabsize=4,
    language=Python,
    inputencoding=utf8,
    extendedchars=true,
    literate=
      {á}{{\'{a}}}1 {é}{{\'{e}}}1 {í}{{\'{\i}}}1 {ó}{{\'{o}}}1 {ú}{{\'{u}}}1
      {ã}{{\~{a}}}1 {õ}{{\~{o}}}1 {ê}{{\^{e}}}1 {ô}{{\^{o}}}1
      {ç}{{\c{c}}}1 {Ç}{{\c{C}}}1
      {Ã}{{\~{A}}}1 {É}{{\'{E}}}1 {Í}{{\'{I}}}1
}

% =====================================================================
% CAPA
% =====================================================================
\begin{document}

{%
\thispagestyle{empty}
\begin{center}

    \includegraphics[width=0.15\textwidth]{figures/logo_urutai.png} \\[0.2cm]
    \vspace{1cm}

    {\small \textbf{INSTITUTO FEDERAL GOIANO} \\}
    {\small \textbf{CAMPUS URUTAÍ} \\}
    {\footnotesize NÚCLEO DE INFORMÁTICA \\}
    {\footnotesize ESPECIALIZAÇÃO EM INTELIGÊNCIA ARTIFICIAL APLICADA A DADOS CORPORATIVOS \\}
    \vspace{2cm}

    {\Large \MakeUppercase{Leonardo Maximino Bernardo} \par}
    \vspace{4.5cm}

    {\Large \textbf{\MakeUppercase{AcolheMente Escolar PB: Motor de Inferência Lógico para Apoio à Gestão Escolar na Priorização Responsável de Acolhimento em Saúde Mental}} \par}

    \vfill

    {\large Urutaí, \today \par}
\end{center}
}
\newpage

% =====================================================================
% SUMÁRIO
% =====================================================================
\pdfbookmark[0]{Sumário}{sumario}
\tableofcontents*
\cleardoublepage

% =====================================================================
% CAPÍTULO 1 — APRESENTAÇÃO DO ALUNO
% =====================================================================
\chapter{Apresentação do Aluno e Contexto}

Leonardo Maximino Bernardo é professor de Física da rede pública estadual da Paraíba desde 2020, com mais de 15 anos de experiência em docência nos níveis médio e superior. Doutor em Ciência e Engenharia de Materiais pela Universidade Federal da Paraíba, sua trajetória acadêmica combina modelagem computacional avançada e uma crescente atuação em ciência de dados e inteligência artificial.

Ao longo de sua formação, atuou como pesquisador na UFPB, desenvolvendo rotinas de cálculo e algoritmos para análise de grandes volumes de dados de simulação numérica de solidificação de ligas metálicas, extraindo padrões quantitativos para publicações acadêmicas. Essa experiência consolidou habilidades de pensamento analítico, resolução de problemas complexos e comunicação científica que informam diretamente sua prática docente e sua abordagem em projetos de inteligência artificial.

Atualmente, além da docência em Física, cursa Pós-Graduação \textit{Lato Sensu} em Inteligência Artificial Aplicada pelo Instituto Federal Goiano e possui formação em Ciência de Dados pelo Unipê. Sua \textit{stack} tecnológica inclui Python, JavaScript, SQL, \textit{frameworks} como Django e React, e ferramentas de análise como Pandas, Scikit-learn e NetworkX. Possui certificações em desenvolvimento \textit{back-end} com Node.js, Harvard CS50 e aceleração global de desenvolvimento de software.

O presente trabalho nasce da intersecção entre sua vivência como professor de escola pública estadual na Paraíba --- onde testemunha diariamente demandas crescentes de acolhimento em saúde mental de adolescentes --- e seu domínio técnico em modelagem computacional e lógica formal. A escolha do domínio de saúde mental escolar reflete o compromisso de aplicar inteligência artificial de forma ética, explicável e governada em um contexto de vulnerabilidade social, onde erros tecnológicos podem ter consequências humanas graves.

% =====================================================================
% CAPÍTULO 2 — INTRODUÇÃO
% =====================================================================
\chapter{Introdução}

A saúde mental de adolescentes no Brasil atravessa um momento crítico. Dados da Pesquisa Nacional de Saúde do Escolar --- PeNSE 2024 \cite{ibge_pense_2024} --- revelam que parcela significativa dos estudantes de 13 a 17 anos reporta sentimentos persistentes de tristeza, solidão e desesperança, com índices particularmente preocupantes na região Nordeste e, especificamente, no estado da Paraíba. O relatório da Organização Mundial da Saúde \cite{who_mental_health_2022} confirma essa tendência global: transtornos mentais respondem por uma proporção crescente da carga de doença entre jovens, e a maioria dos casos permanece sem qualquer forma de acolhimento adequado.

No contexto da escola pública estadual paraibana, a demanda por acolhimento em saúde mental frequentemente supera a capacidade de resposta da equipe pedagógica e psicossocial. Professores, coordenadores e orientadores educacionais enfrentam o desafio de identificar, dentre centenas de estudantes, aqueles que necessitam de atenção prioritária --- sem dispor de ferramentas estruturadas para essa tarefa. O resultado é uma dependência excessiva de percepções individuais, sujeitas a viés cognitivo, fadiga decisória e inconsistência entre turnos e unidades escolares. Essa lacuna não é trivial: a ausência de priorização sistemática pode significar que um adolescente em situação de risco grave não receba acolhimento em tempo hábil.

Ao mesmo tempo, a adoção acrítica de tecnologias de inteligência artificial nesse domínio apresenta riscos igualmente sérios. Sistemas baseados em aprendizado de máquina ou IA generativa, embora poderosos em outros contextos, introduzem opacidade decisória incompatível com o rigor ético exigido quando se trata de menores de idade em situação de vulnerabilidade \cite{jobin_ai_ethics_2019}. A Lei Geral de Proteção de Dados --- LGPD \cite{brasil_lgpd_2018} --- e o Estatuto da Criança e do Adolescente --- ECA \cite{brasil_eca_1990} --- impõem salvaguardas rigorosas para o tratamento de dados sensíveis de menores, incluindo o direito à explicação de decisões automatizadas. Um modelo de ``caixa preta'' que classifique ou ranqueie estudantes por risco psicológico seria, no mínimo, juridicamente temerário e, no limite, eticamente inaceitável.

Diante desse cenário, o presente trabalho propõe o \textbf{AcolheMente Escolar PB}: um Sistema Baseado em Conhecimento que utiliza Lógica Proposicional com Inferência por Resolução para apoiar --- e nunca substituir --- a gestão escolar na priorização responsável de acolhimento em saúde mental. A escolha por lógica simbólica determinística não é uma limitação técnica, mas uma decisão de \textit{design} ético: garante explicabilidade total, rastreabilidade de cada conclusão e ausência de qualquer linguagem diagnóstica clínica \cite{russell_norvig_2022}.

Este documento está organizado da seguinte forma: o Capítulo 3 define formalmente o problema e os objetivos; o Capítulo 4 apresenta a fundamentação metodológica; os Capítulos 5 a 9 detalham dados, representação do conhecimento, regras, conversão para CNF e mecanismo de resolução; o Capítulo 10 descreve a implementação computacional; o Capítulo 11 apresenta resultados; o Capítulo 12 discute explicabilidade; o Capítulo 13 oferece discussão crítica; o Capítulo 14 analisa perspectivas de evolução; e o Capítulo 15 conclui o trabalho.

% =====================================================================
% CAPÍTULO 3 — DEFINIÇÃO DO PROBLEMA E OBJETIVOS
% =====================================================================
\chapter{Definição do Problema e Objetivos}

O problema central que este trabalho endereça pode ser formulado da seguinte maneira: como apoiar a gestão escolar de escolas públicas estaduais da Paraíba na priorização sistemática de acolhimento em saúde mental de adolescentes, utilizando inteligência artificial explicável, sem produzir diagnósticos clínicos, sem violar a LGPD ou o ECA, e sem substituir o julgamento humano profissional?

O escopo da solução proposta compreende a construção de um motor de inferência lógico determinístico com 7 variáveis proposicionais e 6 regras em Forma Normal Conjuntiva, alimentado por indicadores agregados da PeNSE 2024 e validado com cenários sintéticos. O sistema não diagnostica, não classifica alunos individualmente, não utiliza dados identificáveis, não substitui avaliação profissional de psicólogos ou assistentes sociais, e não se conecta a APIs de IA generativa.

O \textbf{objetivo geral} é desenvolver e validar um motor de inferência por resolução, baseado em lógica proposicional e fundamentado em dados oficiais da PeNSE 2024, para apoiar a priorização de acolhimento em saúde mental escolar no estado da Paraíba.

Os \textbf{objetivos específicos} são:

\begin{enumerate}
    \item Modelar o domínio de saúde mental escolar em 7 variáveis proposicionais (5 nucleares e 2 contextuais), ancoradas em variáveis da PeNSE 2024.
    \item Formalizar 6 regras de inferência em Forma Normal Conjuntiva e implementar motor de resolução determinístico.
    \item Validar o motor com cenários sintéticos representativos e comparar com \textit{baseline} manual.
    \item Garantir conformidade integral com LGPD, ECA e Programa Saúde na Escola \cite{brasil_pse_2007}.
    \item Documentar explicabilidade, limitações e perspectivas de evolução com rigor acadêmico.
\end{enumerate}

% =====================================================================
% CAPÍTULO 4 — FUNDAMENTAÇÃO METODOLÓGICA
% =====================================================================
\chapter{Fundamentação Metodológica}

O AcolheMente Escolar PB é, em sua essência, um Sistema Baseado em Conhecimento --- uma classe de sistemas de inteligência artificial que codifica explicitamente o conhecimento de especialistas humanos em regras formais e utiliza mecanismos de inferência para derivar conclusões a partir de fatos observados \cite{brachman_levesque_2004}. Diferentemente de abordagens estatísticas ou conexionistas, um SBC não ``aprende'' a partir de dados brutos: opera sobre uma base de conhecimento construída deliberadamente por engenheiros do conhecimento em colaboração com especialistas do domínio.

Essa característica é particularmente valiosa no contexto de saúde mental escolar, onde o volume de dados individuais disponível é insuficiente para treinamento de modelos de \textit{machine learning}, a transparência de cada decisão é requisito ético e legal, e o custo de um erro --- falso positivo que rotula indevidamente um estudante, ou falso negativo que ignora um caso grave --- é inaceitavelmente alto.

A técnica específica adotada é a Lógica Proposicional com Inferência por Resolução. Trata-se de um formalismo da lógica clássica no qual o conhecimento é representado por proposições atômicas e conectivos lógicos. A inferência é realizada por meio do método de resolução, um procedimento correto e refutacionalmente completo para a lógica proposicional \cite{genesereth_nilsson_1987}. O método opera sobre cláusulas em Forma Normal Conjuntiva: para verificar se uma conclusão $A$ pode ser derivada de uma base de conhecimento $KB$, o procedimento adiciona a negação da conclusão ($\neg A$) ao conjunto de cláusulas e aplica repetidamente a regra de resolução, buscando derivar a cláusula vazia, o que constitui prova por contradição de que $KB \models A$.

A Trilha I do curso de Pós-Graduação em IA Aplicada do IF Goiano tem como foco técnicas simbólicas e de busca. A escolha por lógica proposicional com resolução é duplamente justificada: atende ao escopo pedagógico da trilha e oferece as propriedades de explicabilidade e determinismo exigidas pelo domínio.

Foram consideradas e descartadas alternativas como Lógica de Primeira Ordem, que introduziria complexidade desnecessária para um domínio com 7 variáveis binárias; Redes Bayesianas, que exigiriam distribuições de probabilidade condicional inestimáveis a partir de dados agregados; Árvores de Decisão, que dependem de treinamento supervisionado com dados rotulados inexistentes; Redes Neurais, incompatíveis com os requisitos de explicabilidade total; e IA Generativa, descartada por não oferecer garantias de consistência lógica e determinismo.

Conforme argumentam \citeonline{russell_norvig_2022}, a escolha da representação do conhecimento deve ser guiada pelas propriedades do domínio, e não pela sofisticação da técnica. No caso do AcolheMente Escolar PB, a lógica proposicional é a ferramenta \textit{correta}, não uma ferramenta \textit{limitada}.

O princípio de \textit{human-in-the-loop} é inegociável. A variável de saída $A$ é um sinal para o profissional humano, não uma decisão final. O fluxo operacional previsto é: o motor sinaliza priorização; a equipe pedagógica avalia o caso; a decisão de ação é sempre humana.

% =====================================================================
% CAPÍTULO 5 — DADOS E GOVERNANÇA
% =====================================================================
\chapter{Dados e Governança}

O AcolheMente Escolar PB opera com uma arquitetura rigorosa de proveniência de dados, dividida em três camadas mutuamente exclusivas. Essa separação não é apenas uma decisão técnica: é um requisito de governança que impede contaminação cruzada entre fontes de naturezas distintas.

A primeira camada, denominada TIER\_A, corresponde à PeNSE 2024, base de dados oficial e de acesso público do IBGE \cite{ibge_pense_2024}. A PeNSE é uma pesquisa amostral com representatividade nacional, regional e por unidade da federação. Os microdados utilizados referem-se exclusivamente ao Tema 12 (Saúde Mental) e ao Tema 01 (Informações Gerais), com recorte para o estado da Paraíba. É fundamental ressaltar que a PeNSE não constitui prontuário individual: trata-se de dados agregados, anonimizados na origem, sem possibilidade de identificação de respondentes.

A segunda camada, TIER\_B, compreende cenários sintéticos gerados computacionalmente para demonstração operacional e validação do motor lógico. Cada cenário é um vetor binário de 7 posições, representando combinações hipotéticas de indicadores. Nenhum dado real de aluno ou escola é utilizado, o que permite verificar a correção lógica do motor sem qualquer risco de privacidade.

A terceira camada, TIER\_C, é um \textit{dataset} externo utilizado exclusivamente para \textit{benchmark} metodológico --- para validar que o \textit{pipeline} de processamento funciona corretamente com dados tabulares de outra origem. O \textit{dataset} externo não gera conclusões sobre o Brasil, o Nordeste ou a Paraíba. Qualquer tentativa de inferência regional a partir do TIER\_C é logicamente inválida e explicitamente proibida na arquitetura do sistema.

\textbf{Merge entre tiers é estritamente proibido.} A separação garante que conclusões derivadas de dados oficiais nunca sejam contaminadas por dados sintéticos ou externos, e vice-versa.

A LGPD \cite{brasil_lgpd_2018} classifica dados de saúde de menores como dados pessoais sensíveis, impondo requisitos particularmente rigorosos. O AcolheMente atende a esses requisitos \textit{by design}: não processa dados pessoais identificáveis, opera sobre indicadores agregados anonimizados, toda conclusão é rastreável a regras explícitas, e não produz decisões automatizadas sobre indivíduos.

O ECA \cite{brasil_eca_1990} veda qualquer forma de rotulagem, classificação ou estigmatização de crianças e adolescentes. O motor lógico incorpora \textit{guardrails} que impedem uso de linguagem diagnóstica clínica, \textit{ranking} comparativo entre estudantes e classificação de alunos individuais reais.

O Programa Saúde na Escola \cite{brasil_pse_2007} estabelece a integração entre escola e rede de saúde como política pública. O AcolheMente não substitui essa integração: apoia-a, oferecendo à equipe escolar um instrumento de priorização que indica quando acionar os fluxos de encaminhamento previstos no PSE.

\begin{table}[H]
\centering
\caption{Arquitetura de proveniência de dados em três camadas.}
\label{tab:tiers}
\begin{tabular}{lll}
\toprule
\textbf{Tier} & \textbf{Fonte} & \textbf{Finalidade} \\
\midrule
TIER\_A & PeNSE 2024 (IBGE)    & Inferência oficial para PB \\
TIER\_B & Escola sintética      & Validação do motor \\
TIER\_C & Dataset externo       & Benchmark metodológico \\
\bottomrule
\end{tabular}
\end{table}

% =====================================================================
% CAPÍTULO 6 — REPRESENTAÇÃO DO CONHECIMENTO
% =====================================================================
\chapter{Representação do Conhecimento}

A modelagem do domínio seguiu princípios clássicos de engenharia do conhecimento \cite{brachman_levesque_2004}: identificação de conceitos relevantes, elicitação de regras com especialistas, formalização em lógica proposicional e validação iterativa. O resultado é um conjunto de 7 variáveis proposicionais --- 5 nucleares e 2 contextuais --- que representam o espaço de estados relevante para a priorização de acolhimento.

A decisão de utilizar exatamente 7 variáveis foi deliberada. Com 5 variáveis, perder-se-ia a capacidade de modelar fatores contextuais que modulam a urgência do acolhimento. Com mais de 7, o espaço combinatório de $2^n$ estados cresceria sem ganho proporcional de discriminação, violando o princípio de parcimônia que governa sistemas simbólicos interpretáveis.

As cinco variáveis nucleares são: $E$ (sofrimento emocional recorrente, fontes PeNSE B12004, B12005, B12007), $B$ (baixo apoio socioafetivo percebido, fontes B12003, B07004), $V$ (indicador crítico de desvalor da vida, fonte B12008), $S$ (sinal autorreferido de autoagressão, fonte B12009) e $A$ (acolhimento humano prioritário, variável de saída inferida pelo motor).

As duas variáveis contextuais de modulação são: $C$ (contexto comportamental agravante, fontes B03010C, B03006B) e $I$ (insuficiência institucional de resposta, fontes E01P60, E01P117). As variáveis contextuais nunca inferem $A$ isoladamente. $C$ só participa de regras em conjunção com $E$ e $B$; $I$ só participa em conjunção com $V$. Essa restrição é um \textit{guardrail} lógico que impede falsos positivos por fatores contextuais.

O sistema não diagnostica condições clínicas. As variáveis representam indicadores comportamentais observáveis e auto-reportados, não categorias nosológicas. A variável $A$ sinaliza necessidade de acolhimento pedagógico e psicossocial --- uma ação escolar, não uma intervenção clínica.


\begin{table}[H]
\centering
\caption{Variáveis proposicionais do AcolheMente Escolar PB.}
\label{tab:variaveis}
\begin{tabular}{clll}
\toprule
\textbf{Var} & \textbf{Tipo} & \textbf{Semântica Operacional} & \textbf{Fontes PeNSE} \\
\midrule
E & Nuclear    & Sofrimento emocional recorrente       & B12004, B12005, B12007 \\
B & Nuclear    & Baixo apoio socioafetivo percebido     & B12003, B07004 \\
V & Nuclear    & Indicador de desvalor da vida          & B12008 \\
S & Nuclear    & Sinal de autoagressão                  & B12009 \\
A & Nuclear    & Acolhimento prioritário (saída)        & Inferida \\
C & Contextual & Contexto comportamental agravante      & B03010C, B03006B \\
I & Contextual & Insuficiência institucional             & E01P60, E01P117 \\
\bottomrule
\end{tabular}
\end{table}


% =====================================================================
% CAPÍTULO 7 — BASE DE CONHECIMENTO E REGRAS LÓGICAS
% =====================================================================
\chapter{Base de Conhecimento e Regras Lógicas}

A base de conhecimento do AcolheMente é composta por 6 regras de inferência, cada uma representando um cenário em que a combinação de indicadores justifica a priorização de acolhimento. As regras foram definidas com base em princípios pedagógicos e psicossociais, não em correlações estatísticas.

A regra R1, $S \rightarrow A$, estabelece que a presença de sinal autorreferido de autoagressão é suficiente para acionar incondicionalmente o fluxo de acolhimento. Esta é a regra de maior criticidade: autoagressão é o indicador mais grave do espectro modelado, e qualquer atraso no acolhimento pode ter consequências irreversíveis.

A regra R2, $V \land E \rightarrow A$, captura a coocorrência de desvalor da vida e sofrimento emocional recorrente, indicando um quadro de vulnerabilidade emocional persistente. As regras R3 ($E \land B \rightarrow A$) e R4 ($V \land B \rightarrow A$) modelam cenários em que sofrimento emocional ou desvalor da vida se combinam com baixo apoio socioafetivo, configurando situações em que o estudante não dispõe de rede de suporte.

A regra R5, $E \land B \land C \rightarrow A$, incorpora o contexto comportamental agravante como modulador, enquanto R6, $V \land I \rightarrow A$, combina desvalor da vida com insuficiência institucional de resposta. Note-se que as variáveis contextuais ($C$ e $I$) só participam em conjunção com variáveis nucleares, nunca isoladamente.

A variável $A$ é a única conclusão derivável pelo motor. Essa decisão reflete o princípio de que o sistema deve emitir um sinal binário claro, sem gradações ou categorizações intermediárias que possam ser mal interpretadas pela equipe escolar.


\begin{table}[H]
\centering
\caption{Regras lógicas R1--R6 e respectivas cláusulas CNF.}
\label{tab:regras}
\begin{tabular}{cll}
\toprule
\textbf{Regra} & \textbf{Forma Implicativa} & \textbf{CNF} \\
\midrule
R1 & $S \rightarrow A$              & $\neg S \lor A$ \\
R2 & $V \land E \rightarrow A$      & $\neg V \lor \neg E \lor A$ \\
R3 & $E \land B \rightarrow A$      & $\neg E \lor \neg B \lor A$ \\
R4 & $V \land B \rightarrow A$      & $\neg V \lor \neg B \lor A$ \\
R5 & $E \land B \land C \rightarrow A$ & $\neg E \lor \neg B \lor \neg C \lor A$ \\
R6 & $V \land I \rightarrow A$      & $\neg V \lor \neg I \lor A$ \\
\bottomrule
\end{tabular}
\end{table}


% =====================================================================
% CAPÍTULO 8 — CONVERSÃO PARA CNF
% =====================================================================
\chapter{Conversão para Forma Normal Conjuntiva}

Para aplicar o algoritmo de inferência por resolução, todas as regras implicativas devem ser convertidas para Forma Normal Conjuntiva --- uma conjunção de cláusulas disjuntivas de literais. A transformação segue a equivalência lógica fundamental $p \rightarrow q \equiv \neg p \lor q$. Para regras com antecedentes conjuntivos, a equivalência generaliza-se: $p \land q \rightarrow r \equiv \neg p \lor \neg q \lor r$.

Esse procedimento é correto e preserva a validade lógica das regras originais \cite{russell_norvig_2022}. A CNF resultante é o formato de entrada para o motor de resolução. A Tabela~\ref{tab:regras} apresenta tanto as formas implicativas quanto as cláusulas CNF correspondentes.

A base de conhecimento em CNF possui propriedades desejáveis: é satisfatível, não contém cláusulas contraditórias, $A$ é derivável sempre que ao menos uma regra tem todos os literais de antecedente verdadeiros, e $A$ não é derivável quando apenas variáveis contextuais são verdadeiras, confirmando o \textit{guardrail} de \textit{design}.

A correção da conversão pode ser verificada por tabela-verdade: para cada regra $R_i$, toda valoração que satisfaz a forma implicativa também satisfaz a cláusula CNF correspondente, e vice-versa. Essa equivalência é uma propriedade fundamental da lógica proposicional clássica e independe do domínio de aplicação.

% =====================================================================
% CAPÍTULO 9 — INFERÊNCIA POR RESOLUÇÃO
% =====================================================================
\chapter{Inferência por Resolução}

O método de resolução é um procedimento de prova por refutação: para verificar se $KB \models A$, adiciona-se $\neg A$ ao conjunto de cláusulas e aplica-se iterativamente a regra de resolução. Se a cláusula vazia $\square$ for derivada, conclui-se que $KB \models A$; caso contrário, $A$ não pode ser derivada da base de conhecimento.

A regra de resolução opera sobre pares de cláusulas que contenham literais complementares. Dadas duas cláusulas $C_1 = \{l, \alpha_1, \ldots\}$ e $C_2 = \{\neg l, \beta_1, \ldots\}$, o resolvente é $C_r = \{\alpha_1, \ldots, \beta_1, \ldots\}$. O procedimento é completo para a lógica proposicional: se $KB \models A$, o método sempre encontra a prova em tempo finito.

O pseudocódigo do motor de resolução implementado é:

\begin{lstlisting}[caption={Pseudocódigo do motor de resolução.},label={lst:pseudocodigo}]
funcao RESOLUCAO(KB, conclusao):
    clausulas <- KB uniao {NOT conclusao}
    repita:
        novas <- conjunto vazio
        para cada par (Ci, Cj) em clausulas:
            resolvente <- RESOLVER(Ci, Cj)
            se resolvente = clausula_vazia:
                retorna VERDADEIRO
            novas <- novas uniao {resolvente}
        se novas subconjunto de clausulas:
            retorna FALSO
        clausulas <- clausulas uniao novas
\end{lstlisting}

A implementação completa em Python encontra-se no Apêndice C. O motor foi projetado para ser mínimo e auditável: menos de 100 linhas de código, sem dependências externas além da biblioteca padrão.

A complexidade computacional do método de resolução para lógica proposicional é exponencial no pior caso, mas para a base de conhecimento do AcolheMente (6 cláusulas com no máximo 4 literais cada), a convergência é virtualmente instantânea.

% =====================================================================
% CAPÍTULO 10 — IMPLEMENTAÇÃO COMPUTACIONAL
% =====================================================================
\chapter{Implementação Computacional}

O motor de inferência foi implementado em Python, seguindo princípios de modularidade e testabilidade. A arquitetura separa três responsabilidades: representação das cláusulas CNF (Apêndice B), mecanismo de resolução (Apêndice C) e interface de validação com cenários sintéticos (Apêndice D).

A implementação adota o padrão de representação de cláusulas como conjuntos imutáveis (\texttt{frozenset}) de literais, onde literais positivos são representados pelo nome da variável e literais negativos pelo prefixo \texttt{\textasciitilde}. Essa escolha garante que cláusulas sejam \textit{hashable} e possam ser armazenadas em conjuntos, evitando duplicatas durante o processo de resolução.

O motor inclui função de explicabilidade que, para cada conclusão $A = \text{verdadeiro}$, identifica quais regras foram acionadas e produz texto explicativo em linguagem natural. Essa funcionalidade é essencial para o princípio de \textit{human-in-the-loop}: a equipe escolar pode verificar não apenas \textit{se} o acolhimento foi priorizado, mas \textit{por que}.

A tentativa acadêmica v2 possui 42 testes locais dedicados à validação do contrato abnTeX2, da estrutura acadêmica e da qualidade do PDF gerado. A suíte técnica consolidada do projeto principal --- que verifica correção de cada regra R1--R6, \textit{guardrails} de variáveis contextuais, cenário nulo, cenário completo, explicabilidade, ausência de rótulos diagnósticos e ausência de nós representando estudantes individuais --- permanece separada e não foi alterada nesta tentativa.

O código-fonte completo encontra-se nos Apêndices A--F.

% =====================================================================
% CAPÍTULO 11 — RESULTADOS E COMPARAÇÕES
% =====================================================================
\chapter{Resultados e Comparações}

O motor foi validado com o conjunto exaustivo de $2^6 = 64$ combinações possíveis das variáveis de entrada. Para cada cenário, verificou-se se a conclusão do motor está correta em relação às regras formais, quais regras foram acionadas e se a rastreabilidade é completa. A Tabela~\ref{tab:cenarios} apresenta os cenários representativos de validação.

Os resultados confirmam o comportamento esperado em todos os cenários testados. Quando nenhuma variável é verdadeira (cenário nulo), $A$ não é inferida, demonstrando ausência de falsos positivos em condições de ausência de indicadores. A variável $S$ (autoagressão) é suficiente para acionar o acolhimento, confirmando a prioridade incondicional da regra R1. As variáveis contextuais, quando verdadeiras isoladamente ou em combinação ($C$, $I$, $C+I$), não acionam $A$, confirmando o \textit{guardrail} arquitetural.

A comparação entre o modelo operacional atual (decisão manual pela equipe escolar) e o motor lógico evidencia ganhos em transparência, consistência e explicabilidade. No \textit{baseline} manual, a decisão depende da percepção individual do profissional, que pode variar entre turnos, dias e unidades escolares. O motor lógico oferece 100\% de auditabilidade: cada conclusão é rastreável a regras explícitas em CNF. O \textit{baseline} manual está sujeito a fadiga decisória; o motor é determinístico.

Entretanto, o motor não captura nuances qualitativas que um profissional experiente percebe em interações face a face. Não substitui a escuta ativa, a empatia ou o julgamento clínico. É uma ferramenta de apoio, não de substituição. A existência de um sistema automatizado pode gerar a percepção equivocada de que o acolhimento está ``resolvido'' pela tecnologia --- risco que deve ser mitigado por formação continuada.


\begin{table}[H]
\centering
\caption{Cenários sintéticos de validação do motor de inferência.}
\label{tab:cenarios}
\begin{tabular}{lcccccccl}
\toprule
\textbf{Cenário} & \textbf{E} & \textbf{B} & \textbf{V} & \textbf{S} & \textbf{C} & \textbf{I} & \textbf{A} & \textbf{Regra(s)} \\
\midrule
Nulo             & 0 & 0 & 0 & 0 & 0 & 0 & NÃO & --- \\
S isolado        & 0 & 0 & 0 & 1 & 0 & 0 & SIM  & R1 \\
V+E              & 1 & 0 & 1 & 0 & 0 & 0 & SIM  & R2 \\
E+B              & 1 & 1 & 0 & 0 & 0 & 0 & SIM  & R3 \\
V+B              & 0 & 1 & 1 & 0 & 0 & 0 & SIM  & R4 \\
E+B+C            & 1 & 1 & 0 & 0 & 1 & 0 & SIM  & R3, R5 \\
V+I              & 0 & 0 & 1 & 0 & 0 & 1 & SIM  & R6 \\
C isolado        & 0 & 0 & 0 & 0 & 1 & 0 & NÃO & --- \\
I isolado        & 0 & 0 & 0 & 0 & 0 & 1 & NÃO & --- \\
Todas            & 1 & 1 & 1 & 1 & 1 & 1 & SIM  & R1--R6 \\
\bottomrule
\end{tabular}
\end{table}


\begin{table}[H]
\centering
\caption{Comparação: \textit{baseline} manual vs. motor lógico.}
\label{tab:baseline}
\begin{tabular}{lcc}
\toprule
\textbf{Critério} & \textbf{Manual} & \textbf{Motor Lógico} \\
\midrule
Transparência     & Baixa      & Total \\
Consistência      & Variável   & Determinística \\
Explicabilidade   & Verbal     & Formal (CNF) \\
Rastreabilidade   & Nenhuma    & Regra por regra \\
Fadiga decisória  & Presente   & Ausente \\
Nuance qualitativa & Alta      & Nenhuma \\
\bottomrule
\end{tabular}
\end{table}

% =====================================================================
% CAPÍTULO 12 — EXPLICABILIDADE
% =====================================================================
\chapter{Explicabilidade}

A explicabilidade não é um atributo desejável no AcolheMente Escolar PB: é um requisito ético, jurídico e pedagógico simultaneamente. Do ponto de vista ético, a aplicação de inteligência artificial a populações vulneráveis --- adolescentes em contexto escolar --- exige que cada decisão ou sinal produzido pelo sistema seja compreensível por profissionais não técnicos. Do ponto de vista jurídico, a LGPD \cite{brasil_lgpd_2018} garante ao titular de dados o direito de solicitar explicação sobre decisões tomadas com base em tratamento automatizado. Do ponto de vista pedagógico, um instrumento de apoio à gestão escolar só será adotado se a equipe pedagógica compreender --- e confiar em --- seu funcionamento.

\citeonline{ribeiro_lime_2016} definem explicabilidade como a capacidade de um sistema de apresentar, em termos compreensíveis, as razões que levaram a uma determinada conclusão. Em modelos de \textit{machine learning}, isso frequentemente exige técnicas \textit{post-hoc} como LIME ou SHAP, que aproximam o comportamento de um modelo opaco. No AcolheMente, a explicabilidade é \textbf{intrínseca}: a lógica proposicional é, por definição, transparente. Cada variável tem semântica clara, cada regra tem antecedentes e consequente explícitos, e cada conclusão pode ser rastreada à cláusula CNF que a originou.

O sistema implementa três mecanismos complementares de explicabilidade. O primeiro é a rastreabilidade de regras: para cada conclusão $A = \text{verdadeiro}$, o sistema identifica explicitamente quais regras foram acionadas. O profissional escolar pode verificar, por exemplo, que ``o acolhimento foi priorizado porque a regra R3 foi ativada, ou seja, há sofrimento emocional recorrente e baixo apoio socioafetivo''. O segundo mecanismo é a rastreabilidade de fontes: cada variável proposicional é mapeada a variáveis específicas da PeNSE 2024, permitindo compreender a origem conceitual dos indicadores. O terceiro é o grafo de explicabilidade (Apêndice E), que visualiza as relações entre variáveis de entrada, regras e conclusão.

É importante reconhecer os limites da explicabilidade lógica. O sistema explica \textit{por que} uma conclusão foi derivada (quais regras foram acionadas), mas não explica \textit{por que} as regras são corretas --- essa justificativa reside na engenharia do conhecimento, na literatura de saúde mental escolar e no julgamento dos especialistas que validaram as regras. Além disso, o grafo de explicabilidade pode sugerir relações causais que não existem: a ativação de uma regra indica correlação lógica entre premissas e conclusão, não causalidade clínica.

A relação entre explicabilidade e os marcos regulatórios é direta. O ECA \cite{brasil_eca_1990} exige proteção integral de crianças e adolescentes, o que inclui proteção contra classificação automatizada opaca. A LGPD \cite{brasil_lgpd_2018} garante transparência no tratamento de dados. O PSE \cite{brasil_pse_2007} prevê que ações de saúde na escola sejam realizadas por profissionais capacitados, com compreensão plena dos instrumentos utilizados. A explicabilidade intrínseca do AcolheMente atende simultaneamente a esses três marcos.

A distinção entre explicabilidade intrínseca e \textit{post-hoc} é particularmente relevante no domínio de saúde mental escolar. Enquanto técnicas \textit{post-hoc} geram explicações aproximadas que podem não refletir o verdadeiro comportamento do modelo, a lógica proposicional garante que a explicação \textit{é} o modelo. Não há gap entre o que o sistema faz e o que o sistema diz que faz.

O princípio de \textit{human-in-the-loop} complementa a explicabilidade: mesmo com um sistema totalmente transparente, a decisão final de acolhimento é sempre humana. O motor sinaliza; o profissional decide. Essa separação é inviolável.

% =====================================================================
% CAPÍTULO 13 — DISCUSSÃO CRÍTICA
% =====================================================================
\chapter{Discussão Crítica}

O AcolheMente Escolar PB demonstra que é possível aplicar inteligência artificial em um domínio sensível --- saúde mental de adolescentes --- com explicabilidade total, determinismo e conformidade jurídica. A escolha por lógica proposicional, embora tecnicamente simples quando comparada a redes neurais profundas, é precisamente adequada ao escopo do problema. A arquitetura de governança --- com separação de \textit{tiers}, \textit{guardrails} de não diagnóstico, proibição de \textit{merge} entre fontes e \textit{human-in-the-loop} obrigatório --- representa uma contribuição metodológica que transcende o motor lógico em si.

O sistema apresenta limitações que devem ser explicitamente reconhecidas. A lógica proposicional não permite modelar gradações, incertezas ou relações temporais. Um estudante que apresenta sofrimento emocional ``leve'' e outro com sofrimento ``severo'' são representados pela mesma variável binária $E = \text{verdadeiro}$. A qualidade das conclusões depende inteiramente da qualidade da alimentação das variáveis de entrada: se um profissional atribui $E = \text{falso}$ quando o indicador deveria ser verdadeiro, o motor não pode corrigir o erro.

O motor não aprende com novos dados. Se padrões de vulnerabilidade mudarem ao longo do tempo, as regras precisam ser atualizadas manualmente por engenheiros do conhecimento. A validação foi realizada computacionalmente com cenários sintéticos, não em ambiente escolar real. A transição para uso operacional exigiria piloto institucional com acompanhamento ético.

O principal risco é o falso conforto tecnológico: a existência de um sistema automatizado pode levar a equipe escolar a reduzir a atenção qualitativa aos estudantes, confiando excessivamente no motor. Para mitigar esse risco, a implementação deve ser acompanhada de formação continuada que enfatize o caráter auxiliar --- nunca substitutivo --- do sistema.

Outro risco é o viés de construção: as regras refletem o julgamento dos engenheiros do conhecimento que as definiram. Se esse julgamento contiver vieses implícitos, as regras os reproduzirão. A mitigação exige revisão periódica por comitê multidisciplinar.

A PeNSE 2024 oferece representatividade amostral para a Paraíba, conferindo validade ecológica às variáveis. Entretanto, os indicadores são agregados --- representam tendências populacionais, não perfis individuais. A aplicação a uma escola específica exige calibração contextual.

O AcolheMente não transforma a escola em clínica. O acolhimento escolar é uma ação pedagógica e psicossocial que envolve escuta, encaminhamento e articulação com a rede de proteção social. O motor apoia a priorização dessa ação, não a ação em si. Essa distinção é fundamental para evitar a medicalização do espaço escolar.

A contribuição do projeto reside, portanto, não na sofisticação algorítmica, mas na demonstração de que rigor técnico e responsabilidade ética podem coexistir --- e de que, em domínios de alto risco humano, a transparência é mais valiosa que a complexidade.

% =====================================================================
% CAPÍTULO 14 — PERSPECTIVA DE EVOLUÇÃO DO TRABALHO
% =====================================================================
\chapter{Perspectiva de Evolução do Trabalho}

A arquitetura do AcolheMente Escolar PB foi concebida para permitir evolução incremental, mas essa evolução não é governada pela lógica do ``mais complexo é melhor''. Pelo contrário, qualquer avanço técnico deve ser precedido por validações éticas e institucionais proporcionais ao risco que a mudança introduz. Este capítulo discute os horizontes de evolução possíveis, as pré-condições para cada um e os riscos associados à tentação de substituir o motor simbólico por abordagens estatísticas sem o devido rigor.

O caminho mais natural de evolução é o refinamento do próprio motor simbólico, sem abandonar o paradigma determinístico que lhe confere explicabilidade e rastreabilidade. A incorporação de novas variáveis --- por exemplo, frequência escolar, participação em atividades extracurriculares ou indicadores de convivência familiar --- poderia enriquecer a modelagem sem comprometer a transparência. Essa expansão, entretanto, não é trivial: cada nova variável dobra o espaço combinatório de estados, exigindo revisão integral das regras, revalidação dos cenários sintéticos e reavaliação dos \textit{guardrails}. A engenharia do conhecimento que sustenta o motor atual precisaria ser refeita em colaboração com especialistas do domínio, e não apenas por decisão técnica unilateral.

Uma evolução particularmente promissora dentro do paradigma simbólico é a transição para lógica \textit{fuzzy}, com graus de pertinência entre 0 e 1. Essa abordagem permitiria modelar gradações de sofrimento emocional e apoio social --- por exemplo, distinguir entre um estudante que reporta tristeza ocasional e outro que reporta tristeza persistente --- preservando parte da interpretabilidade da lógica clássica. A definição das funções de pertinência, entretanto, exigiria calibração empírica com dados longitudinais e validação com profissionais de psicologia escolar, o que constitui um projeto de pesquisa em si \cite{russell_norvig_2022}.

Em horizonte mais distante, modelos híbridos poderiam integrar módulos de \textit{machine learning} para tarefas específicas --- como detecção de padrões temporais em séries históricas de indicadores escolares --- mantendo o motor simbólico como camada de governança e explicabilidade. Nessa arquitetura, o componente estatístico operaria sob supervisão estrita do componente simbólico, que atuaria como \textit{guardrail} inviolável: nenhuma conclusão do módulo de \textit{machine learning} seria apresentada à equipe escolar sem passar pelo crivo das regras lógicas. Essa separação arquitetural garante que a opacidade inerente a modelos conexionistas não contamine a interface de decisão humana \cite{jobin_ai_ethics_2019}.

A utilização de redes neurais profundas no domínio de saúde mental escolar permanece uma hipótese estritamente teórica, condicionada ao cumprimento simultâneo de pré-condições rigorosas que, no estado atual da infraestrutura educacional brasileira, não estão atendidas. Do ponto de vista ético, seria necessária aprovação por comitê de ética em pesquisa vinculado à instituição, parecer favorável de especialistas em saúde mental infantojuvenil e avaliação de impacto sobre direitos fundamentais. Do ponto de vista jurídico, a conformidade demonstrável com a LGPD \cite{brasil_lgpd_2018} --- incluindo \textit{Data Protection Impact Assessment} --- e com o ECA \cite{brasil_eca_1990} seria mandatória. Do ponto de vista de dados, seria necessária uma base longitudinal com volume suficiente para treinamento, dados perfeitamente anonimizados e consentidos, e ausência de viés amostral significativo.

Qualquer evolução do sistema, independentemente de sua natureza técnica, deve ser acompanhada de validação institucional robusta. Essa validação inclui piloto em escola-piloto com termo de cooperação técnica, formação continuada da equipe escolar sobre o uso e os limites da ferramenta, criação de comitê de governança de IA com representação multidisciplinar, e avaliação de impacto após período mínimo de uso supervisionado. O \textit{roadmap} de longo prazo prevê integração com os fluxos do Programa Saúde na Escola \cite{brasil_pse_2007}, permitindo que o sinal de priorização alimente formulários de encaminhamento para Unidades Básicas de Saúde e Centros de Atenção Psicossocial --- sem, contudo, produzir dados clínicos ou substituir a avaliação profissional.

Versões futuras devem implementar monitoramento contínuo de viés, abrangendo a distribuição das priorizações por gênero, raça e etnia, localização geográfica e turno escolar. A auditoria humana periódica das regras por comitê multidisciplinar --- com participação de psicólogos, assistentes sociais, pedagogos e representantes da comunidade escolar --- deve ser institucionalizada como prática obrigatória, não opcional. Relatórios de transparência semestrais, com indicadores de uso, taxas de concordância entre o motor e a decisão humana, e \textit{feedback} qualitativo da equipe escolar, constituem instrumentos essenciais de governança. O risco de medicalização da escola --- transformar o espaço pedagógico em extensão do sistema de saúde --- deve ser permanentemente monitorado e mitigado por formação que reforce a natureza pedagógica, e não clínica, do acolhimento escolar.

É necessário resistir à tentação de equiparar sofisticação algorítmica a melhoria do projeto. Aumentar a complexidade do modelo sem ganho proporcional de governança é, no contexto do AcolheMente, um retrocesso. A simplicidade do motor proposicional não é uma limitação a ser superada: é uma \textit{feature} a ser preservada enquanto as condições éticas, jurídicas e institucionais para abordagens mais complexas não estiverem plenamente atendidas. Conforme argumentam \citeonline{russell_norvig_2022}, a escolha da técnica deve ser guiada pelas propriedades do domínio --- e o domínio de saúde mental escolar exige, antes de tudo, confiança, transparência e cautela.

Independentemente da evolução técnica, o compromisso fundamental permanece inalterável: o sistema nunca diagnosticará condições clínicas. Mesmo em versões futuras com componentes de \textit{machine learning}, a camada de governança simbólica deverá atuar como \textit{guardrail} inviolável, garantindo que a saída seja sempre um sinal de priorização --- nunca um rótulo clínico, nunca um \textit{ranking} de estudantes, nunca uma classificação nosológica. Esse compromisso não é uma restrição técnica: é a razão de ser do projeto.

% =====================================================================
% CAPÍTULO 15 — CONCLUSÃO
% =====================================================================
\chapter{Conclusão}

O presente trabalho demonstrou a viabilidade de aplicar inteligência artificial simbólica --- especificamente, lógica proposicional com inferência por resolução --- ao domínio sensível de saúde mental escolar. O AcolheMente Escolar PB, motor de inferência com 7 variáveis proposicionais e 6 regras em Forma Normal Conjuntiva, oferece uma contribuição técnica e metodológica ao debate sobre IA responsável em contextos educacionais.

Do ponto de vista técnico, o trabalho contribui com a formalização de um domínio complexo em lógica proposicional, demonstrando que problemas reais de apoio à decisão podem ser modelados com técnicas da Trilha I sem recurso a \textit{machine learning}; com a implementação de motor de resolução determinístico validado contra todos os $2^6$ cenários possíveis; e com a construção de \textit{pipeline} reprodutível com testes automatizados para validação contínua.

Do ponto de vista metodológico, a contribuição reside na arquitetura de governança de dados em três \textit{tiers}, nos \textit{guardrails} que impedem linguagem diagnóstica e \textit{ranking} de estudantes, na adoção de \textit{human-in-the-loop} como princípio inegociável, e na conformidade demonstrável com LGPD, ECA e PSE.

Do ponto de vista escolar, o trabalho oferece uma ferramenta de apoio à gestão que reduz a dependência de percepções individuais, mecanismo de explicabilidade que permite auditar cada decisão, e integração conceitual com os fluxos do Programa Saúde na Escola.

Os limites são reconhecidos: a expressividade restrita, a dependência de alimentação humana correta, a validação limitada a cenários sintéticos e o risco de falso conforto tecnológico. Esses limites não são falhas de \textit{design}: são consequências deliberadas de um sistema que prioriza segurança, transparência e cautela.

A maturidade de um projeto de IA aplicada a populações vulneráveis não se mede pela complexidade do modelo, mas pela robustez de sua governança. O AcolheMente Escolar PB demonstra que é possível ser tecnicamente rigoroso e eticamente responsável --- e que, em domínios de alto risco humano, essa combinação não é opcional.

% =====================================================================
% REFERÊNCIAS
% =====================================================================
\bibliographystyle{abntex2-alf}
\bibliography{referencias}

% =====================================================================
% APÊNDICES
% =====================================================================
\begin{apendicesenv}
\partapendices

\chapter{Símbolos Proposicionais do AcolheMente}
\label{ap:simbolos}
\lstinputlisting[caption={Símbolos proposicionais do AcolheMente.}, label={lst:simbolos}]{appendices/simbolos_acolhemente.py}

\chapter{Regras Lógicas e CNF}
\label{ap:regras}
\lstinputlisting[caption={Regras lógicas e CNF.}, label={lst:regras}]{appendices/regras_acolhemente.py}

\chapter{Motor de Inferência por Resolução}
\label{ap:motor}
\lstinputlisting[caption={Motor de inferência por resolução.}, label={lst:motor}]{appendices/motor_resolucao_acolhemente.py}

\chapter{Cenários Sintéticos de Validação}
\label{ap:cenarios}
\lstinputlisting[caption={Cenários sintéticos de validação.}, label={lst:cenarios}]{appendices/cenarios_sinteticos_acolhemente.py}

\chapter{Grafo de Explicabilidade}
\label{ap:grafo}
\lstinputlisting[caption={Gerador do grafo de explicabilidade.}, label={lst:grafo}]{appendices/grafo_explicabilidade_acolhemente.py}

\chapter{Gerador de Tabelas de Resultados}
\label{ap:tabelas}
\lstinputlisting[caption={Pipeline de geração de tabelas LaTeX.}, label={lst:tabelas}]{appendices/gerar_tabelas_resultados.py}

\end{apendicesenv}

\end{document}
